\documentclass[12pt,a4paper]{article}
\usepackage{polski}
\usepackage[utf8]{inputenc}
\usepackage{graphicx}
\usepackage[backend=bibtex]{biblatex}
\newcommand{\nazwa}{tekst}

\addbibresource{referencje.bib}

\begin{document}
\begin{titlepage}
\centering
\includegraphics[width=3cm]{logoPWr_EN_pion_czarne.jpg}

\vspace{1cm}
{\Large Politechnika Wrocławska}

\vspace{0.5cm}
Wydział Elektroniki, Fotoniki i Mikrosystemów

\vspace{3cm}
{\LARGE Laboratorium technologii informacyjnych}

\vspace{0.5cm}
Ćwiczenie nr 3

\vfill
Alesky S
\end{titlepage}
\tableofcontents
\newpage

\section{Rysunek}

\begin{figure}[h]

\centering

\includegraphics[width=10cm]{wallhaven-je37wy.jpg}
\caption{Las śćiółka}
\label{fig:rys1}

\end{figure}
Na rysunku \ref{fig:rys1} przedstawiono las \cite{sakdaskd}

\section{Tabelka}

\begin{tabular}{l|c|c|r}

Kot & 5 & 12 & \nazwa \\
Pies & 3 & 222 & \ref{fig:rys1} \\
 sad  & 2 & 222 & 3234 \\
 
\end{tabular}
\section{Równanie}

\begin{equation}
    I=I_S \left( e^{qV/nkT}-1\right)
\end{equation}
$I$ — prąd diody
$I_S$ — prąd nasycenia
$q$ — ładunek elementarny
$V$ — napięcie na złączu
$n$ — wspólczynnik idealności
$k$ — stała Boltzmanna
$T$ — temperature bezwzględną

\printbibliography

\end{document}
